problem:  
Let \((M^{n+1}, g)\) be a globally hyperbolic Lorentzian manifold with compact Cauchy hypersurfaces, satisfying the **timelike convergence condition** \(\mathrm{Ric}(X,X) \ge 0\) for every timelike vector \(X\).  

Suppose there exists a smooth, **spacelike Cauchy hypersurface** \(\Sigma \subset M\) containing a compact, two-sided **stable marginally outer trapped surface** \(S\subset\Sigma\) whose stability operator \(\mathcal{L}\) (in the sense of Andersson–Mars–Simon) has nonnegative principal eigenvalue \(\lambda_1(\mathcal{L}) \ge 0\).  

Assume further that \(\Sigma\) has **nonnegative mean curvature** (with respect to the future-pointing normal) and that outside the domain bounded by \(S\), all outer null expansions are strictly positive.  

**Question:**  
Under these hypotheses, is \(M\) necessarily isometric (locally or globally) to a Lorentzian warped product of the form  
\[
(I\times N, -dt^2 + f(t)^2 h),
\]
where \((N, h)\) is a Riemannian manifold with \(\mathrm{Ric}_h \ge 0\), and the warping function \(f\) satisfies an evolution inequality derived from the Raychaudhuri equation for null congruences starting from \(S\)?  

Equivalently: can the existence of a *stable* marginally outer trapped surface in a nonnegatively curved, globally hyperbolic spacetime enforce a **rigidity-type splitting** analogous to the Cheeger–Gromoll or Hawking–Ellis splitting phenomena, governed by geometric inequalities connecting the null mean curvature and the warping function’s differential equation?  

---

Why is it a "good" problem:  
This problem refines the attempt to link **stability of marginally outer trapped surfaces** (MOTS)—a quasi-local signature of gravitational confinement—with **global rigidity phenomena** in Lorentzian comparison geometry. It seeks a framework where the analytical stability condition of a MOTS translates into an equality or inequality in the Raychaudhuri or Ricci comparison equations, potentially leading to a global classification of spacetime geometry.  

1. **Profound insight:** The problem asks whether a quasi-local geometric property (existence of a stable MOTS) imposes a global warped-product structure, a conceptual analogue of rigidity theorems in Riemannian geometry but through null expansions instead of Ricci tensors of gradient fields.  

2. **Bridge between fields:** It connects **general relativity** (through the study of trapped surfaces and the Raychaudhuri equation) with **comparison and rigidity geometry** in Lorentzian manifolds. It would blend PDE stability analysis, mean curvature evolution, and global causality—a multidomain integration rare in existing theory.  

3. **Extreme simplicity, deep consequences:** In form, the question is almost elementary—“Does one stable MOTS determine the spacetime’s large-scale geometry?”—while its resolution would require combining tools from geometric analysis, energy conditions, and causal geometry in fundamentally new ways.  

4. **Potential to open new directions:** A positive resolution could yield a Lorentzian counterpart of classical rigidity theorems—demonstrating that the geometric presence of quasi-local horizons forces the ambient spacetime into highly constrained warped-product forms. Conversely, counterexamples would sharply delineate the boundaries of Lorentzian comparison geometry and trapped-surface theory.  

This clarity, cross-disciplinary reach, and depth in a conceptually compact statement make it a genuinely “good” problem.